% --- 颜色定义 ---
\definecolor{goldGround}{RGB}{255, 215, 0}    % 金色背景
\definecolor{royalBlue}{RGB}{0, 35, 102}      % 皇室蓝文字
\definecolor{mathBlue}{RGB}{52, 152, 219}    % 杨辉三角用蓝
\definecolor{mathDark}{RGB}{44, 62, 80}      % 深灰色

% --- 1. 为封面设置局部边距 ---
\newgeometry{left=2cm, right=2cm, top=2.5cm, bottom=2.5cm}
\begin{titlepage}
  \centering

  \begin{tikzpicture}[
      remember picture, overlay,
      opacity=0.2, scale=1, rotate=-30,
      shift={($(current page.north west)+(1.5cm,-2cm)$)},
      node distance=0.4cm,
      num/.style={circle, draw, minimum size=0.4cm, font=\sffamily\tiny},
      prime/.style={num, fill=goldGround, draw=goldGround, thick},
      composite/.style={num, fill=gray!10, draw=gray!50, text=gray!80},
      sieve/.style={->, >=Stealth, thick, bend left=25,draw=gray},
      sieveR/.style={->, >=Stealth, thick, bend right=25,draw=gray}
    ]

    % 1. 绘制数字序列 (2-18)
    \foreach \x in {2,...,25} {
      \node[num] (n\x) at (\x*.9, 0) {\x};
    }

    \node[anchor=west] at (0,0){};

    % 2. 模拟欧拉筛过程

    % --- i = 2 ---
    \node[prime] at (n2) {2}; % 2是质数
    % 筛除：2 * 2 = 4
    \draw[sieve] (n2) to (n4);
    \node[composite] at (n4) {4};

    % --- i = 3 ---
    \node[prime] at (n3) {3}; % 3是质数
    % 筛除：3 * 2 = 6
    \draw[sieveR] (n3) to  (n6);
    \node[composite] at (n6) {6};
    % 筛除：3 * 3 = 9
    \draw[sieveR] (n3) to (n9);
    \node[composite] at (n9) {9};

    % --- i = 4 ---
    % 4 已是合数，其最小质因子是2
    % 筛除：4 * 2 = 8
    \draw[sieve] (n4) to  (n8);
    \node[composite] at (n8) {8};
    % 注意：此处 i % p == 0 (4%2==0)，停止筛除，不再筛除 4*3=12

    % --- i = 5 ---
    \node[prime] at (n5) {5};
    % 筛除：5 * 2 = 10
    \draw[sieveR] (n5) to (n10);
    \node[composite] at (n10) {10};
    \draw[sieveR] (n5) to (n15);
    \node[composite] at (n15) {15};
    \draw[sieveR] (n5) to (n25);
    \node[composite] at (n25) {25};

    % --- i = 6 ---
    % 筛除：6 * 2 = 12
    \draw[sieve] (n6) to (n12);
    \node[composite] at (n12) {12};

    % --- i = 7 ---
    \node[prime] at (n7) {7};
    % 筛除：7 * 2 = 14
    \draw[sieveR] (n7) to (n14);
    \node[composite] at (n14) {14};
    \draw[sieveR] (n7) to (n21);
    \node[composite] at (n21) {21};
    
    % --- i = 8 ---
    % 筛除：8 * 2 = 16
    \draw[sieve] (n8) to (n16);
    \node[composite] at (n16) {16};
    
    % --- i = 9 ---
    % 筛除：9 * 2 = 18
    \draw[sieveR] (n9) to (n18);
    \node[composite] at (n18) {18};
    
    % --- i = 10 ---
    % 筛除：10 * 2 = 20
    \draw[sieve] (n10) to (n20);
    \node[composite] at (n20) {20};

    % --- i = 11 ---
    \node[prime] at (n11) {11};
    % 筛除：11 * 2 = 22
    \draw[sieveR] (n11) to (n22);
    \node[composite] at (n22) {22};
    
    % --- i = 12 ---
    % 筛除：12 * 2 = 24
    \draw[sieve] (n12) to (n24);
    \node[composite] at (n24) {24};

    \node[prime] at (n13) {13};
    \node[prime] at (n17) {17};
    \node[prime] at (n19) {19};
    \node[prime] at (n23) {23};

  \end{tikzpicture}


  % --- 顶部几何图案装饰 ---
  \begin{tikzpicture}[
      remember picture,overlay,
      opacity=.1,
      shift={($(current page.north)+(5.5cm,-5cm)$)},
      scale=1.8, rotate=15
    ]
      \draw[royalBlue, very thick] (0,0) -- (3,0) -- (0,4) -- cycle;
      \node at (1.5,-0.3) {$a$};
      \node at (-0.3,2) {$b$};
      \node at (1.7,2.2) {$c$};
  \end{tikzpicture}

  % 核心视觉：空间感杨辉三角 (不使用 scope 嵌套，彻底避免报错) ---
  % 投影原理：通过设置 x 和 y 的倾斜向量来模拟双轴旋转 (Y轴30°, X轴45°)
  \begin{tikzpicture}[
      remember picture, overlay,
      shift={($(6cm,-9cm)$)},
      scale=1.3,rotate=-20
    ]
    \foreach \n in {0,...,10} {

      \foreach \k in {0,...,\n} {
        % 准确计算杨辉三角数值
        %\pgfmathparse{int(round(factorial(\n)/(factorial(\k)*factorial(\n-\k))))}
        \edef\val{\fpeval{trunc(fact(\n)/(fact(\k)*fact(\n-\k)), 0)}}
       
        % 坐标逻辑：在倾斜坐标系下的严谨布局
        % 我们不再使用 \posX \posY，直接在 node 坐标里写 (\k, \n) 
        % 这样 TikZ 会自动按照上面定义的 x, y 倾斜向量进行整体投影，绝不散架

        % 选择性点亮：边缘和中轴
        \pgfmathparse{ (\k == 0 || \k == \n || (\n>0 && \k == \n/2)) ? 1 : 0 }
        \ifdim \pgfmathresult pt > 0.5pt
        \tikzset{cellStyle/.style={draw=mathBlue!10, fill=mathBlue!5, text=mathDark!10, thin}}
        \else
        \tikzset{cellStyle/.style={draw=mathBlue!10, fill=mathBlue!2, text=mathDark!10, thin}}
        \fi

        % 基于 (k, n) 绘制。因为 x 和 y 向量已经倾斜，所以图形整体会呈空间感。
        % (k - n/2) 保证了每一层相对于顶点居中
        \node at ({(\k - \n/2)}, {-\n}) [
          cellStyle,
          opacity=1,%\layerOpacity,
          regular polygon, regular polygon sides=6, 
          minimum size=1cm,
          inner sep=0pt,
          font=\footnotesize\bfseries,
          % 这里的 rotate 补偿了坐标轴倾斜，让六边形看起来还是“平放”的
          rotate=-20
        ] {
          % 这里的 rotate 补偿了坐标轴倾斜，使数字水平
          \rotatebox{-4}{\val}
        };
      }
    }
  \end{tikzpicture}
  % --- 金底蓝字绶带标题 ---
  % 采用 TikZ 直接绘制绶带效果，避免 segmentation 报错
  \begin{tikzpicture}[ baseline,
    yshift={-1*\textheight*.382+1.75cm}]
    % 绶带阴影/折叠部分
    \fill[goldGround!60!black] (-0.2,-0.2) rectangle (0.5, 3.7);
    \fill[goldGround!60!black] (13.5,-0.2) rectangle (14.2, 3.7);

    % 主绶带主体
    \node[
      fill=goldGround, 
      draw=goldGround!70!black,
      minimum width=14cm, 
      minimum height=3.5cm, 
      drop shadow={black!40!white},
      inner sep=0pt
    ] (ribbon) at (7, 1.75) {
      \begin{tabular}{c}
        {\Huge \bfseries \color{royalBlue} 初等数学笔记} \\[0.3cm]
        {\Large \sffamily \color{royalBlue} Elementary Mathematics Notes}
      \end{tabular}
    };
  \end{tikzpicture} 
\end{titlepage}

\restoregeometry

